\documentclass[10pt,a4paper]{amsart}
\usepackage[margin=19mm]{geometry}
\usepackage{amsmath,amssymb,mathtools}
\usepackage[scr=rsfs,cal=euler]{mathalfa}
\usepackage{xfrac}
\usepackage[unicode,hidelinks,bookmarks=false]{hyperref}
\newcommand{\ie}{i.e.\ }
\newcommand{\cf}{cf.\ }
\newcommand{\resp}{respectively\ }
\newcommand{\Subsection}{\subsection}
\providecommand{\nobreakdash}{\nobreak\hbox{-}}
\setlength{\emergencystretch}{2em}
\multlinegap0pt
\usepackage{amssymb}
\usepackage{enumerate}
\usepackage{tikz}
\usepackage{float}
\usepackage{mathtools}
\newtheorem{theorem}{Theorem}
\newtheorem{corollary}{Corollary}
\newtheorem{lemma}{Lemma}
\DeclareMathOperator{\height}{ht}
\newcommand \hs{\operatorname{HS}}
\title{On the Hilbert series and multiplicity of Generalized Binomial Edge Ideals}
\author{Paramhans Kushwaha}
\date{Source: arXiv:2609.12730v1 (CC BY 4.0)}
\begin{document}
\maketitle

\noindent\emph{Source and licence.} On the Hilbert series and multiplicity of Generalized Binomial Edge Ideals. Version arXiv:2609.12730v1 (CC BY 4.0), distributed by arXiv under the Creative Commons Attribution 4.0 International licence, http://creativecommons.org/licenses/by/4.0/, as recorded in the arXiv licence metadata for this exact version and shown on the abstract page https://arxiv.org/abs/2609.12730v1. Full licence notice: \texttt{licenses/arxiv-2609.12730-CC-BY-4.0.txt}.\par\smallskip

\noindent\textit{Editorial scope.} Counted unit: the article's theorem graph (source0.tex lines 583--585). The article's complete original proof of it is delivered with it (source0.tex lines 586--596, one \texttt{proof} environment of 11 lines). No mathematics is written, repaired or re-derived: the statement and the proof are the article's own lines, in the article's own notation, and no retained line is substituted or re-flowed.  Retained uncounted as the source-local context the proof needs: lines 394--400; lines 414--420; lines 556--570.  Nothing was omitted from any retained span, so the package carries no omission record.  No retained line contains a citation, so the wrapper carries no bibliography and no bibliography entry.  No retained line contains a figure environment, an image-inclusion command or drawing code, so no image is delivered and the package declares no asset dependency.  Editorial formatting only, in addition: an AMS \texttt{amsart} preamble that restates the article's own declarations and macros used by the retained lines, namely the environments \texttt{theorem}, \texttt{corollary}, \texttt{lemma} on separate counters, together with the standard packages those lines need.  The retained environments print in this document as Theorem 1, Corollary 1, Lemma 1 in order of appearance, whereas the article prints the same items with the numbering of its own full document.  The complete native source of the article as distributed in the arXiv e-print for this exact version is bundled unchanged as \texttt{sources/210145/source0.tex}.
 \begin{theorem}\label{thm: hs of join of graphs}
     Let $H$  and $H'$ be two graphs on $[p]$ and $[q]$, respectively. Let $G=H*H'$ be the join of $H$ and $H'$. Then 
     \begin{align*}
        \hs(S/J_{m,H*H'},t)=&\hs(S_H/J_{m,H},t)+\hs(S_{H'}/J_{m,H'},t)\\
        & +\frac{\sum_{i\geq 0}\binom{m-1}{i}\left[\binom{p+q-1}{i}-\binom{p-1}{i}(1-t)^q-\binom{q-1}{i}(1-t)^p\right]t^i}{(1-t)^{m+p+q-1}}.
    \end{align*}
 \end{theorem}

 \begin{corollary}\label{cor: dim of join of graphs}
    Let $H$ and $H'$ be two graphs on $[p]$ and $[q]$, respectively. Then 
    \begin{enumerate}
        \item $\dim(S/J_{m,H*H'})=\max\{\dim(S_H/J_{m,H}),\dim(S_{H'}/J_{m,H'}),m+p+q-1\}$.
        \item $\height(J_{m.H*H'})=\min\{mq+\height(J_{m,H}),mp+\height(J_{m,H'}),(m-1)(p+q-1)\}$.
    \end{enumerate}
\end{corollary} 

\begin{lemma}\label{lemma: dimension and multiplicity of path graphs}
    Let $G=P_n$ be the path graph on $n$ vertices. Then
    \begin{enumerate}
        \item \label{lemma: dimension of path graphs} $$ \dim(S/J_{m,G})=
    \begin{cases}
        \lceil\frac{n}{2}\rceil m+1, \text{ if $n$ is even},\\
        \lceil\frac{n}{2}\rceil m,\quad\quad\text{ if $n$ is odd}.
    \end{cases}$$
    \item \label{lemma: multiplicity of path graphs} $$ e_0(S/J_{m,G})=
    \begin{cases}
        \dfrac{mn}{2}, \text{ if $n$ is even},\\
       1,\quad\text{  if $n$ is odd}.
    \end{cases}$$
    \end{enumerate}
\end{lemma}

\begin{theorem}\label{thm: hs of multifan graph}
    Let $G=M_{n_1,\ldots,n_r}$, $r\geq 2$ be a multifan graph and $n=\sum_{i=1}^rn_i$. Then $$\dim(S/J_{m,G})=\sum_{i=1}^r\dim(S/J_{m,P_{n_i}})\text{ and }e_0(S/J_{m,G})=\prod_{n_i=\text{ even }}\dfrac{mn_i}{2}.$$
\end{theorem}

\begin{proof}
It follows from Corollary \ref{cor: dim of join of graphs} that $\dim(S/J_{m,G})=\max\{\dim(S/J_{m,\sqcup_{i=1}^rP_{n_i}}),m+n\}$. Let $d=\dim(S/J_{m,\sqcup_{i=1}^rP_{n_i}})$. Then, we have $d=\sum_{i=1}^r\dim(S/J_{m,P_{n_i}})$.
Let $A=\{n_i:n_i\text{ is even }\}$, $B=\{n_i:n_i\text{ is odd }\}$. Using Lemma \ref{lemma: dimension and multiplicity of path graphs}\eqref{lemma: dimension of path graphs}, we obtain $$d=\sum_{n_i\in A}(\frac{mn_i}{2}+1)+\sum_{n_i\in B}\frac{m(n_i+1)}{2}.$$ Now, we show that $d-m-n>0$. Note that $d-m-n=\frac{1}{2}(m(n-2)-2n+2|A|+m|B|)$. Since $m\geq 3$ and $n\geq 2|A|+|B|$, we get $d-m-n\geq 4|A|+(m+1)|B|-6>0$. Thus, we obtain $\dim(S/J_{m,G})=\sum_{i=1}^r\dim(S/J_{m,P_{n_i}})$. 

Next, we compute the multiplicity. Using Theorem \ref{thm: hs of join of graphs}, we get
   \begin{align*}
      \hs(S/J_{m,G})&=\hs(S/J_{m,\sqcup_{i=1}^rP_{n_i}},t)+\frac{\sum_{i\geq 0}\binom{m-1}{i}\left[\binom{n}{i}-(1-t)\binom{n-1}{i}\right]t^i}{(1-t)^{m+n}}\\
      &=\frac{Q(t)}{(1-t)^{d}}+\frac{\sum_{i\geq 0}\binom{m-1}{i}\left[\binom{n}{i}-(1-t)\binom{n-1}{i}\right]t^i}{(1-t)^{m+n}},
   \end{align*}
    where $Q(t)\in \mathbb Z[t]$. Since $d>m+n$, it follows that  $$e_0(S/J_{m,G})=e_0(S/J_{m,\sqcup_{i=1}^rP_{n_i}})=\prod_{i=1}^r e_0(S/J_{m,P_{n_i}}).$$ By Lemma \ref{lemma: dimension and multiplicity of path graphs}\eqref{lemma: multiplicity of path graphs}, we obtain $e_0(S/J_{m,G})=\prod_{\substack{n_i=\text{ even}}} \dfrac{mn_i}{2}$. This completes the proof.  
\end{proof}

\end{document}
